\documentclass[11pt]{article}
\usepackage[margin=0.65in]{geometry}
\usepackage{booktabs}
\usepackage{float}
\usepackage[T1]{fontenc}
\usepackage[utf8]{inputenc}
\begin{document}
\begin{table}[H]
\centering
\caption{Effect of a 10 Percent Shock to Crude Oil and Gas Prices, 2023 (effects in pp)}
\scriptsize
\setlength{\tabcolsep}{4.2pt}
\begin{tabular}{lrrrrrr}
\toprule
Country & \shortstack{Average inflation\\Total} & \shortstack{Average inflation\\Direct} & \shortstack{Average inflation\\Indirect} & \shortstack{Income inflation\\gap} & \shortstack{Age inflation\\gap} & \shortstack{Rural-urban\\inflation gap}
 \\
\midrule
DE & 0.6 & 0.1 & 0.5 & 0.1 & 0.0 & 0.2
 \\
FR & 0.6 & 0.0 & 0.6 & -0.0 & 0.1 & 0.2
 \\
IT & 1.0 & 0.4 & 0.7 & -0.1 & -0.1 & 0.2
 \\
ES & 0.8 & 0.0 & 0.8 & 0.1 & 0.1 & 0.1
 \\
NL & 1.1 & 0.2 & 0.8 & -0.1 & 0.2 & 0.2
 \\
BE & 0.8 & 0.0 & 0.8 & -0.1 & 0.1 & 0.2
 \\
EA20 & 0.8 & 0.1 & 0.6 & 0.0 & 0.0 & 0.2
 \\
\bottomrule
\end{tabular}
\par\vspace{0.35em}\footnotesize\noindent Sources: Eurostat HICP, HBS, FIGARO, author's calculations.
\par\vspace{0.25em}\footnotesize\noindent Notes. Entries are percentage-point effects of a 10 percent sectoral input price shock. The price response is computed as $\Delta p = (I - A^{\prime})^{-1}s$, with $A$ built from FIGARO. The direct effect is the component from the shocked sector in the household consumption basket. Indirect is the Leontief propagation component. Total = Direct + Indirect before rounding. Income inflation gap = Q1 - Q5 of revenues; age inflation gap = over 60 years - less than 30 years; rural-urban inflation gap = rural areas - cities.
\end{table}
\vspace{0.6em}
\end{document}
